% ============================================================
% f2014_shanghai_arts_q07_views.tex
% 2014 上海卷（文科）第 7 题：割去两个小长方体后的三视图
% 来源 PDF：普通高考/2014/2014上海文.pdf
% 原页：第 1 页（A3），印面 917，填空题第 7 题题干下方
% 原长方体 5（长）×3（深）×4（高）；中间高 2 的一层左右两侧
% 各割去一个 2×2×3 的小长方体，正中留下宽 1 的竖条。
% 正视图：宽 5、高 4 的“工”字形；侧视图：3×4 矩形，高 1、3 处
% 各有实线；俯视图：5×3 矩形，距左边 2、3 处各有一条贯通虚线。
% ============================================================
\begin{tikzpicture}[>=Stealth, line width=0.65pt, font=\small, scale=0.86,
  line cap=round, line join=round,
  dim/.style={<->, >=Stealth, thin, line width=0.45pt}]

  % ---------------- 正视图 ----------------
  \draw (0,0) -- (5,0) -- (5,1) -- (3,1) -- (3,3) -- (5,3) -- (5,4)
        -- (0,4) -- (0,3) -- (2,3) -- (2,1) -- (0,1) -- cycle;

  % 中间竖条宽 1 的尺寸线（延长线伸入上横带）
  \draw[thin] (2,3) -- (2,3.78);
  \draw[thin] (3,3) -- (3,3.78);
  \draw[->, thin, line width=0.45pt] (1.55,3.55) -- (1.96,3.55);
  \draw[->, thin, line width=0.45pt] (3.45,3.55) -- (3.04,3.55);
  \node at (2.5,3.55) {$1$};

  % 左侧三段高度 1、2、1
  \draw[dim] (-0.45,0) -- (-0.45,1);
  \draw[dim] (-0.45,1) -- (-0.45,3);
  \draw[dim] (-0.45,3) -- (-0.45,4);
  \foreach \y in {0,1,3,4} \draw[thin] (-0.62,\y) -- (-0.28,\y);
  \node[left] at (-0.62,0.5) {$1$};
  \node[left] at (-0.62,2) {$2$};
  \node[left] at (-0.62,3.5) {$1$};

  % 底边宽 5
  \draw[dim] (0,-0.5) -- (5,-0.5);
  \draw[thin] (0,-0.3) -- (0,-0.7);
  \draw[thin] (5,-0.3) -- (5,-0.7);
  \node[below] at (2.5,-0.5) {$5$};

  % ---------------- 侧视图 ----------------
  \draw (7,0) rectangle (10,4);
  \draw (7,1) -- (10,1);
  \draw (7,3) -- (10,3);
  \draw[dim] (7,-0.5) -- (10,-0.5);
  \draw[thin] (7,-0.3) -- (7,-0.7);
  \draw[thin] (10,-0.3) -- (10,-0.7);
  \node[below] at (8.5,-0.5) {$3$};

  % ---------------- 俯视图 ----------------
  \draw (0,-6) rectangle (5,-3);
  \draw[dashed] (2,-6) -- (2,-3);
  \draw[dashed] (3,-6) -- (3,-3);

  % 视图名称
  \node at (2.5,4.55) {正视图};
  \node at (8.5,4.55) {侧视图};
  \node at (2.5,-6.55) {俯视图};
\end{tikzpicture}
